\chapter{Einleitung}
Die vorliegende Arbeit beschäftigt sich mit der Konzeption und Umsetzung einer modernen Softwarearchitektur 
für die internen Schnittstellen eines Jahresabschlusssystems im Gesamtkontext einer Enterprise-Steuerberatungssoftware. Diese Schnittstellen sollen primär berechnete bzw. gebuchte Werte aus dem Jahresabschluss 
an die verschiedenen internen Verwendersysteme übergeben. Das gesamte Enterprise-System befindet sich aktuell in einer Transformationsphase von On-Premises-Monolithen hin zu einer dezentralen Cloud-Architektur.
Dieses Kapitel gibt einen Überblick über die Problematik und stellt die Zielsetzung sowie den Aufbau der Arbeit vor.

\section{Problemstellung und Motivation}
Steuerberater müssen sich immer stärker dem Druck der Digitalisierung und Automatisierung stellen. Anforderungen, die ein Steuerberater
2010 an eine Software hatte, sind heute nicht mehr ausreichend. Für Steuerberater ist es daher entscheidend, ihren Betrieb weiter zu digitalisieren, 
um wettbewerbsfähig zu bleiben (\cite{Siegmann2023}, S.~11f.). Innovative Unternehmen suchen nach agilen und digitalisierten Steuerberatern(\cite{Siegmann2023}, S.~10). 
Es ist wichtig für Steuerberater, neue Formen der Kommunikation und des Austauschs mit den Mandanten zu nutzen, die sich durch die zunehmende Digitalisierung ergeben. Hieraus können sich Chancen für eine stärkere Mandantenbindung ergeben (\cite{Siegmann2023}, S.~25).
Darüber hinaus verschieben sich die Aufgaben von Steuerberatern: Wiederholende Tätigkeiten können zunehmend automatisiert werden, was dazu führt, dass Steuerberater und auch deren Angestellte viel mehr analysierende und kontrollierende Tätigkeiten
übernehmen müssen (\cite{Siegmann2023}, S.~36).  Äußere Faktoren haben auch Einfluss auf die Betriebe der Steuerberater. Der Fachkräftemagel gilt hier als eine der größten Herausforderungen für die Branche (hier noch ifo Stats raussuchen).
Durch ortsunabhängiges Arbeiten hoffen Steuerberater, gutes Personal zu finden und langfristig an den Betrieb zu binden (\cite{Siegmann2023}, S.~52). Darüber hinaus besteht die Hoffnung der Betriebe, neue Mandanten durch ortsunabhängige Dienstleistungen zu gewinnen (\cite{Siegmann2023}, S.~51). 

On-Premises-Software kann diese Anforderungen nicht mehr ausreichend erfüllen. Daher ist ein Paradigmenwechsel hin zu Cloud-basierten Softwarelösungen notwendig.
Cloud-Software bietet die Möglichkeit, schneller auf sich ändernde Anforderungen zu reagieren, da sie flexibler und skalierbarer ist als traditionelle On-Premises-Lösungen \cite{Sharma2022}.
So ist On-Premises-Software durch ihre physische Hardware, die beim Steuerberater steht, begrenzt. Dies macht Erweiterungen zeitaufwendig und teuer. In der Cloud kann entsprechend der Nachfrage nach den Ressourcen flexibel skaliert werden. Während Ausfälle oder Fehler in reinen On-Premises-Systemen zu Ausfällen der Systeme führen, können solche Ausfälle in der Cloud durch ihren redundanten Aufbau besser aufgefangen werden. So ist es möglich, dem Kunden resilientere und zuverlässigere Systeme zur Verfügung zu stellen. Auch die Auslieferung von Änderungen ist schneller und kosteneffizienter möglich. Das erlaubt es, Innovationen oder Zukunftstechnologien viel schneller in bestehende Workflows zu integrieren \cite[S. 479 f.]{Sharma2022}.

Das einfache Verschieben bestehender Systeme auf Cloud-Infrastrukturen reicht nicht aus, um die genannten Vorteile und Anforderungen zu erreichen.
Hohpe beschrieb bereits 2019, dass für Unternehmen, die eine ausschließliche Infrastruktur-Migration durchführen, die Gefahr besteht, eine sogenannte „Enterprise Non-Cloud“ zu schaffen. Er bezieht sich hier auf das Risiko, dass Unternehmen ihr gesamtes On-Premises-Vorgehen mitsamt der für das Unternehmen bewährten Technologien und Architekturen in ein neues Rechenzentrum überführen. Dies führt weiterhin dazu, dass Potenziale, die mit der Cloud-Migration geschaffen werden sollen, nicht oder nicht vollständig genutzt werden können (vgl. \cite{Hohpe2020}).
Legacy-On-Premises-Systeme basieren oft auf stark gekoppelten Monolithen. Damit die Vorteile von Cloudanwendungen vollumfänglich genutzt werden können, müssen diese Strukturen in Cloud-Konzepte überführt werden. Eine einfache Migration der Systeme ist nicht ausreichend.
Im Rahmen solch komplexer Migrationen muss sichergestellt werden, dass der laufende Geschäftsbetrieb während der Entwicklung und Migration der Cloudsysteme nicht gestört wird. Das bedeutet, dass die Systeme für eine gewisse Zeit in einem Parallelbetrieb laufen müssen. Weiterhin müssen die Systeme untereinander auch kommunizieren können, damit die On-Premises-Software Stück für Stück abgelöst werden kann. Hier müssen Schnittstellen und Erweiterungen in ein Legacy-System gebaut werden, für die es niemals gedacht war.

Daraus lässt sich ableiten, dass für den vollen Mehrwert von Cloudanwendungen eine neue grundlegende Konzeption der Softwarearchitektur notwendig ist. 
Das Auflösen des großen Rechnungswesenmonolithen in fachliche (Mikro-)Services erfordert eine hohe Qualität und Sorgfalt in der Ausarbeitung.


\section{Zielsetzung und Aufbau der Arbeit}
